\input{../tex-inputs/latex-leseansicht-vorspann}

\section[Arthur Schnitzler an Gustav Schwarzkopf, 18. 5. 1923]{L04495 Arthur Schnitzler an Gustav Schwarzkopf, 18. 5. 1923}
\nopagebreak\mylabel{L04495v}
\rehead{ }\normalsize\beginnumbering\briefempfaengerindex{Schwarzkopf, Gustav@\textsc{Schwarzkopf, Gustav}!zzzSchnitzler, Arthur@\emph{von Arthur Schnitzler}!1923-05-181@{18. 5. 1923}|(be}
\toendnotes[C]{\smallbreak\pagebreak[2]}
\correspDesc{Versand  durch Arthur Schnitzler am 18. 5. 1923 in Stockholm
\newline{}Erhalt  durch Gustav Schwarzkopf im Zeitraum [19. 5. 1923
                  – 23. 5. 1923?] in Wien}\toendnotes[C]{\smallbreak}
\Standort{DLA, A:Schnitzler, HS.1985.1.1897.}
\physDesc{Bildpostkarte, 678 Zeichen
\newline{}Handschrift: Bleistift, lateinische Kurrent
\newline{}Versand: Stempel: »\nobreak{}\oindex{Stockholm@\textbf{Stockholm}|pwk}Stockholm 13, {[}18.{]} 5. {[}23{]}\nobreak{}«.  }\toendnotes[C]{\smallbreak}\pstart{}{\pb}Herrn\pend{}\pstart{}Gustav Schwarzkopf\pend{}\pstart{}Wien I\oindex{I., Innere Stadt@\textbf{I., Innere Stadt}|pw}\pend{}\pstart{}Tiefer Graben 17\oindex{Wien@\textbf{Wien}!I., Innere Stadt@\textbf{I., Innere Stadt}!Tiefer Graben 17@\textbf{Tiefer Graben 17}|pw}\pend{}{\bigskip}
\pstart
           \noindent{}\centering{}{\pb}\textcolor{gray}{\textbf{Stockholm, Riddaholmen\oindex{Riddarholmen@\textbf{Riddarholmen}|pw}}}\pend
           \vspace{1em}
\pstart
           \raggedleft{}{\pb}Stockholm\oindex{Stockholm@\textbf{Stockholm}|pw}{ }1\substVorne{}\textsuperscript{7}\substDazwischen{}\textcolor{gray}{8}\substHinten{}. 5. 92\textcolor{gray}{3}\pend
           \vspace{0.5em}
\pstart
           lieber Gustav, so hat mich meine Reise doch auch hieher geführt. In
                  Kopenhagen\oindex{Kopenhagen@\textbf{Kopenhagen}|pw} hatt ich \label{K_L04495-1v}\edtext{2 Vorlesungen\eventindex{Politikens Hus@\textbf{Politikens Hus}!Lesung von Das Schicksal des Freiherrn von Leisenbohg, Weihnachts-Einkäufe, Der grüne Kakadu, 14.5.1923@Lesung von Das Schicksal des Freiherrn von Leisenbohg, Weihnachts-Einkäufe, Der grüne Kakadu, 14.5.1923|pwv}\eventindex{Kungliga Musikaliska Akademien@\textbf{Kungliga Musikaliska Akademien}!Lesung von Dreifache Warnung, Die letzten Masken, Das Schicksal des Freiherrn von Leisenbohg, Weihnachts-Einkäufe, 18.5.1923@Lesung von Dreifache Warnung, Die letzten Masken, Das Schicksal des Freiherrn von Leisenbohg, Weihnachts-Einkäufe, 18.5.1923|pwv}}{\lemma{\textnormal{\emph{2 Vorlesungen}}}\Cendnote{\textnormal{Lesung von Das Schicksal des Freiherrn von Leisenbohg, Weihnachts-Einkäufe, Der grüne Kakadu, 14.5.1923. In: A. S.: \emph{Kulturveranstaltungen}, \url{https://schnitzler-kultur.acdh.oeaw.ac.at/pmb62185.html} und Lesung von Dreifache Warnung, Die letzten Masken, Das Schicksal des Freiherrn von Leisenbohg, Weihnachts-Einkäufe, 18.5.1923. In: A. S.: \emph{Kulturveranstaltungen}, \url{https://schnitzler-kultur.acdh.oeaw.ac.at/pmb62192.html}.}}}\label{K_L04495-1}, u sah \label{K_L04495-2v}\edtext{die große Szene\pwindex{Schnitzler, Arthur 15. 5. 1862 Wien – 21. 10. 1931 ebd.@\textsc{Schnitzler, Arthur} (15. 5. 1862 Wien – 21. 10. 1931 ebd.)!Große Szene@\strich\emph{Große Szene}|pw}\eventindex{Det Kongelige Teater@\textbf{Det Kongelige Teater}!Aufführung von Große Szene und Napoli, 15.5.1923@Aufführung von Große Szene und Napoli, 15.5.1923|pwv}}{\lemma{\textnormal{\emph{die große Szene}}}\Cendnote{\textnormal{Aufführung von Große Szene und Napoli, 15.5.1923. In: A. S.: \emph{Kulturveranstaltungen}, \url{https://schnitzler-kultur.acdh.oeaw.ac.at/pmb58640.html}.}}}\label{K_L04495-2}{ }\introOben{}\uline{ungekürzt}\introOben{}, hier \label{K_L04495-3v}\edtext{les\eventindex{Kungliga Musikaliska Akademien@\textbf{Kungliga Musikaliska Akademien}!Lesung von Dreifache Warnung, Die letzten Masken, Das Schicksal des Freiherrn von Leisenbohg, Weihnachts-Einkäufe, 18.5.1923@Lesung von Dreifache Warnung, Die letzten Masken, Das Schicksal des Freiherrn von Leisenbohg, Weihnachts-Einkäufe, 18.5.1923|pwv} ich
               heut Abend}{\lemma{\textnormal{\emph{les ich
               heut Abend}}}\Cendnote{\textnormal{Lesung von Dreifache Warnung, Die letzten Masken, Das Schicksal des Freiherrn von Leisenbohg, Weihnachts-Einkäufe, 18.5.1923. In: A. S.: \emph{Kulturveranstaltungen}, \url{https://schnitzler-kultur.acdh.oeaw.ac.at/pmb62192.html}.}}}\label{K_L04495-3}, am 22. ist \label{K_L04495-4v}\edtext{die Liebelei\pwindex{Schnitzler, Arthur 15. 5. 1862 Wien – 21. 10. 1931 ebd.@\textsc{Schnitzler, Arthur} (15. 5. 1862 Wien – 21. 10. 1931 ebd.)!Liebelei. Schauspiel in drei Akten@\strich\emph{Liebelei. Schauspiel in drei Akten}|pw}\eventindex{Kungliga Dramatiska Teatern@\textbf{Kungliga Dramatiska Teatern}!Aufführung von Liebelei, 22.5.1923@Aufführung von Liebelei, 22.5.1923|pwv}}{\lemma{\textnormal{\emph{die Liebelei}}}\Cendnote{\textnormal{Aufführung von Liebelei, 22.5.1923. In: A. S.: \emph{Kulturveranstaltungen}, \url{https://schnitzler-kultur.acdh.oeaw.ac.at/pmb62229.html}.}}}\label{K_L04495-4}.\pend
           
\pstart
           – \label{K_L04495-5v}\edtext{Interviewt\pwindex{Bønnelycke, Emil 21.\,3.\,1893 Aarhus – 27.\,11.\,1953 Söndrum@\textsc{Bønnelycke, Emil} (21.\,3.\,1893 Aarhus – 27.\,11.\,1953 Söndrum)!Dichter und Dramatiker Arthur Schnitzler in Kopenhagen@\strich\emph{Der Dichter und Dramatiker Arthur Schnitzler in Kopenhagen}|pwv}\pwindex{Rechendorff, Henrik @\textsc{Rechendorff, Henrik}!En Samtale med Arthur Schnitzler@\strich\emph{En Samtale med Arthur Schnitzler}|pwv}\pwindex{Vogel-Jörgensen, Torkild 29.\,6.\,1891 Svendborg – 4.\,10.\,1972 Hellerup@\textsc{Vogel-Jörgensen, Torkild} (29.\,6.\,1891 Svendborg – 4.\,10.\,1972 Hellerup)!Arthur Schnitzler. En Samtale med en beremt Wiener [1923]@\strich\emph{Arthur Schnitzler. En Samtale med en beremt Wiener [1923]}|pwv}\pwindex{Falkenfleth, Haagen 4.\,2.\,1875 Kopenhagen – 31.\,12.\,1957 ebd.@\textsc{Falkenfleth, Haagen} (4.\,2.\,1875 Kopenhagen – 31.\,12.\,1957 ebd.)!Arthur Schnitzler i København. En Samtale med den berømte østerrigkse Forfatter@\strich\emph{Arthur Schnitzler i København. En Samtale med den berømte østerrigkse Forfatter}|pwv}\pwindex{Archi @\textsc{Archi}!Arthur Schnitzler nu i Stockholm@\strich\emph{Arthur Schnitzler nu i Stockholm}|pwv}\pwindex{Arthur Schnitzler om »Pandoras ask«@\emph{Arthur Schnitzler om »Pandoras ask«}|pwv}\pwindex{Valentin, Guido 17.\,12.\,1895 Göteborg – 9.\,10.\,1952 Mjölby@\textsc{Valentin, Guido} (17.\,12.\,1895 Göteborg – 9.\,10.\,1952 Mjölby)!Anatols diktare på besök i Stockholm@\strich\emph{Anatols diktare på besök i Stockholm}|pwv}\pwindex{Arthur Schnitzler i Stockholm@\emph{Arthur Schnitzler i Stockholm}|pwv}\pwindex{O. N–ll @\textsc{O. N–ll}!Arthur Schnitzler som charmör och som kåsör@\strich\emph{Arthur Schnitzler som charmör och som kåsör}|pwv}\pwindex{Sir Guy @\textsc{Sir Guy}!Arthur Schnitzler, Anatols författare, besöker Stockholm@\strich\emph{Arthur Schnitzler, Anatols författare, besöker Stockholm}|pwv}\pwindex{Arthur Schnitzler lämnade Stockholm i går@\emph{Arthur Schnitzler lämnade Stockholm i går}|pwv}\pwindex{Även Stockholm är bra@\emph{Även Stockholm är bra}|pwv} bisher (genau) von 15 Personen\pwindex{Bønnelycke, Emil 21.\,3.\,1893 Aarhus – 27.\,11.\,1953 Söndrum@\textsc{Bønnelycke, Emil} (21.\,3.\,1893 Aarhus – 27.\,11.\,1953 Söndrum)|pwv}\pwindex{Rechendorff, Henrik @\textsc{Rechendorff, Henrik}|pwv}\pwindex{Vogel-Jörgensen, Torkild 29.\,6.\,1891 Svendborg – 4.\,10.\,1972 Hellerup@\textsc{Vogel-Jörgensen, Torkild} (29.\,6.\,1891 Svendborg – 4.\,10.\,1972 Hellerup)|pwv}\pwindex{Falkenfleth, Haagen 4.\,2.\,1875 Kopenhagen – 31.\,12.\,1957 ebd.@\textsc{Falkenfleth, Haagen} (4.\,2.\,1875 Kopenhagen – 31.\,12.\,1957 ebd.)|pwv}\pwindex{Archi @\textsc{Archi}|pwv}\pwindex{Valentin, Guido 17.\,12.\,1895 Göteborg – 9.\,10.\,1952 Mjölby@\textsc{Valentin, Guido} (17.\,12.\,1895 Göteborg – 9.\,10.\,1952 Mjölby)|pwv}\pwindex{Andersson, Anna Lisa 3.\,6.\,1873 Stockholm – 10.\,3.\,1958 ebd.@\textsc{Andersson, Anna Lisa} (3.\,6.\,1873 Stockholm – 10.\,3.\,1958 ebd.)|pwv}\pwindex{O. N–ll @\textsc{O. N–ll}|pwv}\pwindex{Sir Guy @\textsc{Sir Guy}|pwv}}{\lemma{\textnormal{\emph{Interviewt … Personen}}}\Cendnote{\textnormal{Es
                     sind zwölf Interviews aus Kopenhagen\oindex{Kopenhagen@\textbf{Kopenhagen}|pwk} und Stockholm\oindex{Stockholm@\textbf{Stockholm}|pwk} dokumentiert, siehe Arthur Schnitzler: \emph{»Das
                        Zeitlose ist von kürzester Dauer.« Interviews, Meinungen und Proteste
                        (1891–1931)}. Hg. Martin Anton Müller,
                     Göttingen: \emph{Wallstein}{ }2023, Bd. 1, Nr. 40–51,
                     S. 157–219.}}}\label{K_L04495-5}, – aber da ich nicht verstehe, was sie
               schreiben, ist die Sache irrelevant. Die ganze »Aufwartung« {\pb}amerikanisch\oindex{Vereinigte Staaten von Amerika [USA]@\textbf{Vereinigte Staaten von Amerika [USA]}|pw} – beim Aussteigen aus dem Zug war
               man photographirt, gefilmt u.s.w. Ich befinde mich trotzdem wohl – die Städte,
               Menschen, Landschaft, Atmosphäre thun mir wohl. Leider ist es immer kalt – der
               Frühling ist viele Wochen hinter unserm zurück. – Seien Sie, ebenso Ihr Bruder\pwindex{Schwarzkopf, Max 12.\,6.\,1857 Wien – 14.\,4.\,1928 ebd.@\textsc{Schwarzkopf, Max} (12.\,6.\,1857 Wien – 14.\,4.\,1928 ebd.)|pwv} herzlich gegrüßt\pend
           \pstart Ihr \spacefill\mbox{A.}\pend{}\selectlanguage{ngerman}\endnumbering\briefempfaengerindex{Schwarzkopf, Gustav@\textsc{Schwarzkopf, Gustav}!zzzSchnitzler, Arthur@\emph{von Arthur Schnitzler}!1923-05-181@{18. 5. 1923}|)be}\mylabel{L04495h}
\begin{anhang}
\end{anhang}\newcommand{\dateiname}{L04495}\newcommand{\titel}{Arthur Schnitzler an Gustav Schwarzkopf, 18. 5. 1923}\newcommand{\editorInnen}{Herausgegeben von Jahnke, SelmaMüller, Martin Anton}\input{../tex-inputs/latex-leseansicht-abspann}
